\section{Shape-changing Driven Locomotion and Morphological Computation}
\label{sec:shape-changing}

If compliant material gives a robot the capacity to deform, geometry and topology decide what that deformation can do. This section moves from the question of substrate to the question of shape, and to a body of work in which a robot's form is treated not as a passive scaffold for its actuators but as part of how it computes its behavior.

\subsection{Definition}
\label{subsec:mc-definition}

The idea has a name. \textit{Morphological computation} describes the phenomenon by which a robot's body, through its geometry, material distribution, and mechanical compliance, performs functions that would otherwise need to be specified by a controller. The concept traces back to Pfeifer and Bongard's argument that intelligence does not reside solely in the brain but emerges from the body's interaction with its environment \cite{pfeiferHowBodyShapes2007}, and is developed through a series of subsequent reviews that classify the different ways a body can carry computational load and survey the embodied-computing implementations built on the principle \cite{paul_morphological_2006, mullerWhatMorphologicalComputation2017, hauserMorphologicalComputationPresent2024}.

Tibbits states the design implication most expansively. As he writes, \inlinequote{The notion of morphological computing could be applicable more generally to say that materials and their geometric configurations---how they are put together, with different material properties, or different mechanical and geometric structures---can form the basis of material computing.}\cite{tibbitsThingsFallTogether2021} Two emphases in this formulation matter for the thesis. First, \textit{configuration}, how things are put together, is named alongside material as the basis for computation. Second, the framing is general enough to apply not only to single-body soft robots but to any system in which structure carries computational load.

A practical illustration of the principle is balance-free design. A radially symmetric robot does not need an inertial sensor and a balancing algorithm to remain upright: the geometry itself ensures that no orientation is privileged, and any face can become the ground. A control function that would otherwise be implemented by gyroscope feedback is replaced by a structural choice. This is morphological computation in its most concrete engineering form, and it is exactly the kind of substitution that motivates the rest of the section.

\subsection{Related Works}
\label{subsec:mc-related}

The most direct precedents for the present work approach shape-driven locomotion through the deformation of an integral, polyhedral body. Usevitch et al.'s untethered isoperimetric soft robot \cite{usevitchUntetheredIsoperimetricSoft2020} is built as an octahedral truss whose edges are isoperimetric soft tubes (\textit{\cref{fig:usevitch_isoperimetric}}); air is transferred between tubes so that one perimeter shrinks as another grows while total tube length is conserved, allowing the robot to reshape its octahedron and roll. Nozaki et al.'s Mochibot \cite{nozaki_continuous_2018} pushes the polyhedral idea further: a spherical robot built around a thirty-two-legged rhombic-triacontahedron skeleton achieves continuous shape-changing locomotion through coordinated leg extension, with the high vertex count of the underlying Catalan solid making the continuity possible. In both works, the polyhedral topology determines which deformations are kinematically possible and which gaits are reachable, but each paper studies a single polyhedron in depth rather than treating polyhedra as a comparable family. Wang, Fei, and Liu's FifoBots \cite{wang_fifobots_2019} take a non-polyhedral route to a related idea: foldable soft segments rotate the body about a momentary edge of contact, producing flipping locomotion whose gait is a consequence of segment geometry rather than a controller-specified command.

A second thread distributes shape-change across modular units rather than concentrating it in an integral body. Chin et al.'s AuxBot \cite{chinFlipperStyleLocomotionStrong2023} uses volume-changing modules whose auxetic shells expand by roughly 274\% in 0.8 seconds; each module is mechanically simple and can only expand or contract, but coordinated phase-shifting across multiple connected modules produces flipper-style locomotion of the whole assembly. AuxBot is in many ways a precursor to the questions taken up in the next section, since the complexity of its locomotion arises from how the modules are arranged and synchronized rather than from any single module's behavior. In a similar spirit, Belke et al.'s Mori3 \cite{belke_morphological_2023} assembles triangular units through a physical polygon-meshing mechanism, allowing the same set of modules to be reconfigured for driving, rolling, or walking; and Yoder et al.'s HEXEL \cite{yoderHexagonalElectrohydraulicModules2024} uses hexagonal electrohydraulic modules whose tessellation supports rapid reconfiguration, with the authors noting explicitly that topology and configuration choices were crucial to the robot's mobility. Each of these systems makes geometry visibly matter, but each makes a single geometric commitment and develops it in depth.

\subsection{Gap}
\label{subsec:mc-gap}

What unites these systems is that each chooses a particular geometric or topological primitive (an octahedron, a triacontahedron, a polygon, a hexagonal tessellation, or an auxetic shell) and develops a robot from it. Within each choice, geometry is doing real work: it sets which deformations are possible, which gaits are available, and how much the controller has to specify. The literature, in other words, contains many examples of geometry carrying computational load.

What it does not yet contain is comparison across geometries. None of the works above asks the question that follows naturally once morphological computation is taken seriously as a design principle: how does the choice of geometry itself shape what the robot can do? When the same actuation strategy and the same material are placed into different polyhedral configurations, do the resulting robots differ qualitatively in mobility, in controllability, or in actuator economy? Each existing system picks one geometry and develops it in depth; none studies the geometry as a controlled variable across a family of related designs. The shape is doing computational work, but the question of \textit{which} shape, and \textit{why}, has been left almost entirely to the designer's intuition rather than to systematic comparison.

This is one of the gaps the thesis works in. It is also only half of the picture, however. Each of the systems above is an integral, single-body robot whose geometry was chosen at the design stage and fixed thereafter. The next section turns to a parallel literature in which the body is decomposed into discrete, comparable units that can be rearranged: modular reconfigurable robotics. There, geometry returns in a different form, as the arrangement of repeatable modules rather than the shape of an integral body.
